\chapter{Family-connection bridge и exact-one BdG kernel}

Предыдущий gate показал, что антисимметричный family selector нельзя
реализовать четырёхмерным scalar Majorana Yukawa term. Однако такой оператор
естественно возникает из ковариантной производной вещественного
\(SO(3)\)-triplet. Это даёт локальный bulk--core bridge без ручного
boundary mass.

\section{Допустимый tubular operator}

Пусть family-blind vortex operator на поперечном диске имеет вид
\(D_\perp\) и обладает единственной нормированной вещественной zero mode
\(\psi_0\), отделённой от остального спектра gap \(\Delta_\perp>0\).
Три поколения до включения family connection образуют
\begin{equation}
 \ker(D_\perp\otimes I_3)
 =\operatorname{span}\{\psi_0\}\otimes\mathbb R^3.
\end{equation}

На трубке \(N(\gamma)\simeq D^2\times S^1_L\) введём вещественную
\(SO(3)\)-связность
\begin{equation}
 \mathcal A_s=\frac1L K_n,
 \qquad
 K_n=\theta_*[n]_\times,
 \qquad
 \operatorname{Hol}_\gamma(\mathcal A)=e^{K_n}=R_n(\theta_*).
\end{equation}
Она входит не как scalar bilinear, а через разрешённую covariant derivative
triplet Majorana field:
\begin{equation}
 S_F=\frac12\int_{N(\gamma)}
 \Psi^TC\left[
 D_{a,\Phi}\otimes I_3+
 \Gamma^s\otimes(\partial_s+\mathcal A_s)
 \right]\Psi.
\end{equation}
Поскольку \(K_n^T=-K_n\), это стандартная реальная gauge coupling. После
проекции на \(\psi_0\) она даёт
\begin{equation}
 P_0S_FP_0
 =\frac12\int_\gamma
 \chi^T(\partial_s+L^{-1}K_n)\chi,
\end{equation}
то есть прежний антисимметричный core operator теперь следует из bulk
covariant derivative.

\section{Точный separable spectrum}

Для product metric, постоянной longitudinal connection, family-blind
vortex profile и ранее выведенной полной periodic longitudinal branch
квадрат BdG operator разделяется:
\begin{equation}
 \mathcal D^2
 =D_\perp^2\otimes I_3+
 I\otimes\left[-i\partial_s+\frac{iK_n}{L}\right]^2.
\end{equation}
Так как
\begin{equation}
 \operatorname{Spec}(iK_n)=\{-\theta_*,0,+\theta_*\},
\end{equation}
полный спектр квадратов имеет вид
\begin{equation}
 \lambda_{j,k,\sigma}^2
 =\varepsilon_j^2+
 \left(\frac{2\pi k+\sigma\theta_*}{L}\right)^2,
 \qquad
 k\in\mathbb Z,
 \quad
 \sigma\in\{-1,0,+1\}.
\end{equation}
Здесь \(\varepsilon_0=0\), а \(|\varepsilon_j|\geq\Delta_\perp\) для
\(j\ne0\). Поскольку
\begin{equation}
 0<\theta_*=\arccos\!\left(\frac{26-9\sqrt{15}}{11}\right)<\pi,
\end{equation}
нулевое значение возможно только при
\begin{equation}
 j=0,
 \qquad k=0,
 \qquad \sigma=0.
\end{equation}

\begin{theorem}[Exact-one separable kernel]
В separable tubular model с unit transverse Majorana index и family Wilson
holonomy \(R_n(\theta_*)\) полный BdG operator \(\mathcal D\) имеет ровно одну вещественную
zero mode:
\begin{equation}
 \dim\ker\mathcal D=1.
\end{equation}
Она равна \(\psi_0\otimes n\). Gap над ней ограничен снизу величиной
\begin{equation}
 \delta_{\rm sep}
 =\min\left\{\Delta_\perp,\frac{\theta_*}{L}\right\}.
\end{equation}
Для систолического core \(L=\pi\) family contribution равен
\begin{equation}
 \frac{\theta_*}{\pi}=0.7979250094\ldots.
\end{equation}
\end{theorem}

Это локальная теорема на tubular neighborhood. Она не утверждает, что
постоянная связность с данным углом продолжается на весь
\(\mathbb{RP}^3\) как плоская связность.

Таким образом, в этом классе моделей не требуется ни scalar order parameter
\(\Sigma\), ни отдельный core Yukawa coefficient \(g\): selector является
longitudinal Wilson splitting того же real triplet bundle.

\section{Устойчивость вне точного product ansatz}

Пусть \(\mathcal D_E=\mathcal D+E\), где \(E\) сохраняет class-D reality,
тот же mod-two defect sector и имеет operator norm
\begin{equation}
 \|E\|<\frac12\delta_{\rm sep}.
\end{equation}
Тогда spectral projection на интервал
\((-\delta_{\rm sep}/2,\delta_{\rm sep}/2)\) сохраняет rank one. Particle--hole
symmetry не позволяет единственной непарной eigenvalue уйти от нуля, поэтому
exact-one kernel устойчив, а новый gap не меньше
\begin{equation}
 \delta_E\geq\delta_{\rm sep}-\|E\|>\frac12\delta_{\rm sep}.
\end{equation}

Для perturbation, контролируемой блоками, достаточно более локального
Feshbach-критерия. Если \(P\) проектирует на три исходные core modes,
\(Q=I-P\),
\begin{equation}
 q=\|QEQ\|<\Delta_\perp,
 \qquad
 b=\|PEQ\|,
\end{equation}
то две поперечные family modes остаются обратимыми при
\begin{equation}
 \frac{\theta_*}{L}>
 \frac{b^2}{\Delta_\perp-q}.
\end{equation}
Осевая mode остаётся точной, если perturbation ковариантна относительно той
же connection и annihilates \(n\).

\section{Что закрыто и что осталось}

Закрыты два локальных разрыва:
\begin{enumerate}
 \item антисимметричный core operator получен из допустимой bulk gauge
 coupling, а не из исчезающего scalar bilinear;
 \item для separable tube доказан full-spectrum exact-one kernel, а для
 малых class-D perturbations дана явная область устойчивости.
\end{enumerate}

\section{Торсионное препятствие для плоского ambient extension}

Систолическая геодезическая представляет элемент порядка два в
\begin{equation}
 \pi_1(\mathbb{RP}^3\times S^1)=\mathbb Z_2\times\mathbb Z.
\end{equation}
Поэтому holonomy любой плоской ambient connection обязана удовлетворять
\begin{equation}
 W_\gamma^2=I.
\end{equation}
В \(SO(3)\) это допускает только углы \(0\) и \(\pi\), тогда как
\(0<\theta_*<\pi\) и \(\theta_*\ne\pi\). Следовательно, tubular connection
\(K_n/L\) не имеет плоского ambient extension вдоль систолического core.

\begin{nogocriterion}[Flat torsion Wilson obstruction]
Нельзя одновременно считать \(V_{\rm gap}\) потенциалом плоской ambient
Wilson-моды и помещать её holonomy \(R_n(\theta_*)\) на торсионный
систолический цикл. Такое чтение несовместимо с fundamental-group relation.
\end{nogocriterion}

\section{Два согласованных продолжения}

Первый маршрут использует неторсионный внешний фактор \(S^1\). Тогда defect
в полном carrier читается как \(\gamma\times S^1\): систолическая
\(\gamma\) несёт vortex topology и periodic longitudinal zero mode, а
family Wilson connection живёт на дополнительном круге длины
\(L_1=2\pi R_1\). Separable spectrum становится
\begin{equation}
 \lambda_{j,\ell,k,\sigma}^2
 =\varepsilon_j^2+
 \left(\frac{2\pi\ell}{L_\gamma}\right)^2+
 \left(\frac{2\pi k+\sigma\theta_*}{L_1}\right)^2.
\end{equation}
При periodic total spin branch на дополнительном круге единственный ноль
снова имеет индексы \((j,\ell,k,\sigma)=(0,0,0,0)\), а
\begin{equation}
 \delta_{S^1}
 =\min\left\{
 \Delta_\perp,
 \frac{2\pi}{L_\gamma},
 \frac{\theta_*}{L_1}
 \right\}.
\end{equation}
Для \(L_\gamma=\pi\), \(R_1=1\) family gap равен
\begin{equation}
 \frac{\theta_*}{2\pi}=a_*=0.3989625047\ldots.
\end{equation}
Однако anti-periodic spin branch на внешнем \(S^1\) уничтожает осевую
нулевую моду. Поэтому этот маршрут требует periodic defect sector либо
дополнительной компенсирующей holonomy, выведенной до данных.

\section{Компенсированная spin-ветвь из charge-two Higgs sector}

Компенсация не требует нового непрерывного поля. Пусть геометрическая
spin-структура на внешнем \(S^1\) anti-periodic и \(\eta=1/2\). Для
fermion нечётного \(B-L\)-заряда \(q\) при flat gauge holonomy
\(\alpha=\oint_{S^1}a\) полный momentum равен
\begin{equation}
 p_{k,\sigma}
 =\frac{2\pi(k+\eta)+q\alpha+\sigma\theta_*}{L_1}.
\end{equation}
После конденсации charge-two field gauge group содержит residual
\(\mathbb Z_2\)-holonomy
\begin{equation}
 \alpha=0\quad\text{или}\quad\alpha=\pi.
\end{equation}
Для нечётного \(q\) ветвь \(\alpha=\pi\) сокращает anti-periodic spin sign:
\begin{equation}
 e^{2\pi i\eta}e^{iq\alpha}=(-1)(-1)=+1.
\end{equation}
При этом charge-two condensate однозначен, поскольку
\begin{equation}
 e^{2i\alpha}=+1.
\end{equation}
После целочисленной перенумерации \(k\) спектр возвращается к periodic
формуле, и осевая \(\sigma=0\) zero mode восстанавливается. Следовательно,
compensated branch совместима с уже введённым \(B-L\) field content и не
вводит непрерывного fit-параметра.

Локальное charge-two Higgs action, однако, не различает две residual flat
ветви \(\alpha=0\) и \(\alpha=\pi\). Поэтому закрыта доступность требуемой
spin\(^{G}\)-структуры, но не её динамический выбор. Для полного вывода нужен
topological term, boundary condition или determinant sign, выбирающий
\(\alpha=\pi\) до обращения к neutrino data.

Действительно, для любого минимального локального функционала, зависящего от
\(F_a\), \(|D_a\Phi|^2\) и \(V(|\Phi|)\), две конфигурации
\begin{equation}
 (a,\Phi)=(0,v)
 \quad\text{и}\quad
 (a,\Phi)=\left(\frac{\pi}{L_1}ds,
 ve^{2\pi is/L_1}\right)
\end{equation}
имеют \(F_a=0\), \(D_a\Phi=0\) и одинаковый potential energy. Они не
связаны допустимым gauge transformation при наличии unit-charge fermions,
поскольку требуемое преобразование меняет знак после обхода круга.

\begin{nogocriterion}[Residual-holonomy degeneracy]
Минимальный local charge-two Higgs action не способен динамически выбрать
compensating branch \(\alpha=\pi\) вместо \(\alpha=0\). Выбор требует
nonlocal Wilson functional, topological coupling или квантового вклада
odd-charge sector. Простое объявление condensate не закрывает spin gate.
\end{nogocriterion}

\section{Однопетлевой Hosotani sign gate}

Для massive mode с \(\rho>0\) введём конечную holonomy-dependent часть
\begin{equation}
 I_\rho(\beta)=
 \log\frac{\cosh(2\pi\rho)-\cos(2\pi\beta)}
 {\cosh(2\pi\rho)-1}.
\end{equation}
На фундаментальном интервале \(0\leq\beta\leq1/2\) она монотонно возрастает,
поскольку
\begin{equation}
 \frac{dI_\rho}{d\beta}
 =\frac{2\pi\sin(2\pi\beta)}
 {\cosh(2\pi\rho)-\cos(2\pi\beta)}\geq0.
\end{equation}
Bosonic mode даёт \(+dI_\rho/2\), fermionic mode --- \(-dI_\rho/2\).

Для anti-periodic odd-charge fermion переход \(\alpha:0\to\pi\) меняет
\(\beta:1/2\to0\), поэтому
\begin{equation}
 \Delta V_f
 =V_f(\pi)-V_f(0)
 =\frac d2 I_\rho(1/2)>0.
\end{equation}
Для anti-periodic fermion с \(q=1/3\) конечная фаза равна
\(\beta=2/3\sim1/3\), и
\begin{equation}
 \Delta V_f
 =\frac d2\left[I_\rho(1/2)-I_\rho(1/3)\right]>0.
\end{equation}
Для periodic boson с \(q=1/3\) правильный shift равен \(q/2=1/6\), а не
\(1/3\), но знак остаётся тем же:
\begin{equation}
 \Delta V_b=\frac d2I_\rho(1/6)>0.
\end{equation}
Charge-two boson не различает ветви. Следовательно, минимальный набор
anti-periodic fermions и periodic bosons действительно выбирает
\(\alpha=0\), а compensated branch \(\alpha=\pi\) является более высокой
stationary branch.

При \(\rho\gg1\) правильная асимптотика равна
\begin{equation}
 I_\rho(\beta)
 =2\left[1-\cos(2\pi\beta)\right]e^{-2\pi\rho}
 +O(e^{-4\pi\rho}).
\end{equation}
В частности,
\begin{equation}
 I_\rho(1/2)\sim4e^{-2\pi\rho},
 \qquad
 I_\rho(1/3)\sim3e^{-2\pi\rho},
 \qquad
 I_\rho(1/6)\sim e^{-2\pi\rho}.
\end{equation}
Ранее предложенные коэффициенты \(2\) и \(1/2\) были занижены вдвое;
отношение \(I(1/2):I(1/6)=4:1\) при этом сохраняется.

\section{Geometric versus total periodicity}

Sterile \(N^c\) в compensated construction имеет geometric spin shift
\(\eta=1/2\). При \(\alpha=\pi\) её \emph{total} phase становится
periodic, \(\beta=0\), но Hosotani determinant сравнивает две ветви при
фиксированной geometric \(\eta\):
\begin{equation}
 \beta(0)=\frac12,
 \qquad
 \beta(\pi)=0,
 \qquad
 \Delta V_{N^c}=+\frac d2I_\rho(1/2)>0.
\end{equation}
Следовательно, такая \(N^c\) усиливает выбор \(\alpha=0\), а не создаёт
требуемый отрицательный вклад. Периодичность defect mode вдоль
систолической \(\gamma\), полученная сокращением spin- и torsion-знаков, также
не меняет geometric spin structure на независимом внешнем \(S^1\).

Отрицательный вклад
\(-dI_\rho(1/2)/2\) требует поля с geometric \(\eta=0\), для которого
\(\beta:0\to1/2\). Такого поля существующая compensated \(N^c\)-ветвь не
предоставляет.

Для минимального \(B-L\)-контента одного поколения имеются двенадцать
anti-periodic Weyl components с \(|q|=1/3\) и четыре с нечётным целым
зарядом, включая \(N^c\). При общей \(rho\) их вклад равен
\begin{equation}
 \Delta V_{gen}
 =2I_\rho(1/2)+6\left[I_\rho(1/2)-I_\rho(1/3)\right]
 =8I_\rho(1/2)-6I_\rho(1/3)>0.
\end{equation}
Даже если искусственно перевести одну \(N^c\) на geometric periodic branch,
получается
\begin{equation}
 \Delta V_{gen}^{\rm retwist}
 =7I_\rho(1/2)-6I_\rho(1/3)>0.
\end{equation}
В large-mass limit эти коэффициенты пропорциональны соответственно
\(14e^{-2\pi\rho}\) и \(10e^{-2\pi\rho}\). Поэтому три retwisted \(N^c\)
при одинаковых masses не проходят magnitude gate.

Для существующего минимального контента common-\(\rho\) assumption не
нужна. При произвольных положительных \(\rho_i\) каждый отдельный вклад имеет
положительный знак:
\begin{align}
 \Delta V_{odd,i}
 &=\frac{d_i}{2}I_{\rho_i}(1/2)>0,\\
 \Delta V_{1/3,i}
 &=\frac{d_i}{2}\left[
 I_{\rho_i}(1/2)-I_{\rho_i}(1/3)
 \right]>0,\\
 \Delta V_{b,1/3,i}
 &=\frac{d_i}{2}I_{\rho_i}(1/6)>0.
\end{align}
Charge-two bosons дают ноль. Следовательно, сумма положительна при любом
mass spectrum существующих anti-periodic fermions и periodic bosons.

\begin{theorem}[Minimal-content multi-mass Hosotani no-go]
При фиксированных geometric spin structures минимальный \(B-L\)-контент
выбирает \(\alpha=0\) для любых \(\rho_i>0\). Полный mass spectrum не может
изменить знак, поскольку отрицательных summands в этом field content нет.
Multi-mass gate остаётся открытым только после добавления нового
geometric-periodic odd-charge sector.
\end{theorem}

\begin{nogocriterion}[Total-periodicity sign no-go]
Нельзя использовать zero mode, которая становится total-periodic только в
ветви \(\alpha=\pi\), как periodic-spin determinant, стабилизирующий эту же
ветвь. Это меняет fixed boundary problem между концами сравнения и обращает
знак Hosotani contribution. Для отрицательного вклада нужен независимо
выведенный geometric-periodic odd-charge sector.
\end{nogocriterion}

\section{Почему обычный CS/theta term не стабилизирует ветвь}

Квантованный topological term различает сектора на уровне фазы. В евклидовой
сигнатуре
\begin{equation}
 S_{CS,E}=\frac{ik}{4\pi}\int a\wedge da,
 \qquad
 S_{\theta,E}=\frac{i\theta}{8\pi^2}\int F\wedge F.
\end{equation}
На полностью flat configuration оба локальных density равны нулю. При
добавлении фиксированного magnetic flux \(m\) Chern--Simons term может дать
\begin{equation}
 \Delta S_{CS,E}=ikm\alpha,
\end{equation}
то есть относительную фазу unit modulus, но не отрицательный вклад в
\(\operatorname{Re}V_{eff}\). Его real Hessian по \(\alpha\) равен нулю.
Discrete theta term аналогично меняет знак или фазу веса topological sector,
но сам по себе не превращает Hosotani maximum в энергетический minimum.

\begin{nogocriterion}[Topological-phase stabilization no-go]
Обычный квантованный CS/theta term без дополнительного constraint sector не
может стабилизировать \(\alpha=\pi\) против положительного real Hosotani
splitting. Он способен изменить interference секторов или, в расширенной
TQFT, спроецировать часть Hilbert space, но это не эквивалентно минимуму
вещественного effective potential.
\end{nogocriterion}

Для реального выбора \(\alpha=\pi\) требуется дополнительный вклад
\begin{equation}
 \Delta V_{min}:=V_{min}(\pi)-V_{min}(0)>0,
 \qquad
 \Delta V_{new}(\pi)-\Delta V_{new}(0)<-\Delta V_{min}.
\end{equation}
Минимальный знаково подходящий пример --- periodic odd-charge fermion: он
даёт
\begin{equation}
 \Delta V_{per\,f}=-\frac d2I_\rho(1/2)<0.
\end{equation}
Его multiplicity и mass должны быть выведены из parent action и превысить
положительный вклад исходного spectrum. Альтернатива --- явная TQFT
projection, удаляющая \(\alpha=0\) sector, но тогда это новый дискретный
сектор модели, а не Hosotani stabilization.

\section{Минимальный geometric-periodic selector sector}

Минимальный anomaly-safe кандидат состоит из двух left-handed singlet Weyl
fields
\begin{equation}
 X_+:(B-L)=+1,
 \qquad
 X_-:(B-L)=-1.
\end{equation}
Они образуют vectorlike Dirac pair. Чтобы при общей anti-periodic spin
structure получить effective geometric periodicity только для этой пары,
введём flat \(\mathbb Z_4^X\)-line с charges
\begin{equation}
 z(X_+)=+1,
 \qquad
 z(X_-)=-1\pmod4,
\end{equation}
и holonomy \(-1\) на внешнем \(S^1\). Тогда spin sign и \(\mathbb Z_4^X\)
sign сокращаются, а B--L holonomy переводит
\begin{equation}
 \beta:0\longrightarrow\frac12
 \qquad (\alpha:0\longrightarrow\pi).
\end{equation}
Пара даёт требуемый real contribution
\begin{equation}
 \Delta V_X=-I_{\rho_X}(1/2)<0.
\end{equation}

Непрерывные anomaly sums сокращаются попарно:
\begin{equation}
 \sum q=\sum q^3=0.
\end{equation}
То же происходит для elementary mixed discrete sums при представлении
\(z=\pm1\):
\begin{equation}
 \sum z=\sum z^3=\sum zq^2=\sum z^2q=0.
\end{equation}
Разрешённая симметрийная Dirac mass непрерывно переводит изолированную пару
в тривиально gapped phase без нарушения \(B-L\) или \(\mathbb Z_4^X\).
Следовательно, pair не несёт остаточной discrete 't Hooft anomaly: anomaly
gate для самого нового fermionic content закрыт.

Более того, Dirac mass \(M_XX_+X_-\) разрешена, тогда как
\begin{equation}
 \Phi X_+X_+,
 \qquad
 \Phi^\dagger X_-X_-,
 \qquad
 N^cX_-
\end{equation}
имеют ненулевой \(\mathbb Z_4^X\)-charge и запрещены. Поэтому новая пара не
участвует в vortex pairing operator и не добавляет defect zero modes;
исходный exact-one kernel сохраняется.

Пусть для benchmark все три минимальных поколения имеют \(\rho=1\). Их
положительный вклад равен
\begin{equation}
 \Delta V_{min}^{(3)}
 =3\left[8I_1(1/2)-6I_1(1/3)\right]
 =0.0783386447\ldots.
\end{equation}
Одна vectorlike pair выбирает \(\alpha=\pi\), если
\begin{equation}
 I_{\rho_X}(1/2)>\Delta V_{min}^{(3)},
\end{equation}
что даёт
\begin{equation}
 \rho_X<0.6259781084\ldots.
\end{equation}
При одинаковой массе \(\rho_X=1\) требуется не менее одиннадцати таких
pairs. В выбранной \(\alpha=\pi\)-ветви даже bare-light pair получает
Wilson gap; при одной паре на threshold
\begin{equation}
 m_{phys}R_1\geq
 \sqrt{\rho_X^2+\frac14}=0.801\ldots,
\end{equation}
поэтому determinant leverage не обязательно создаёт massless spectator в
физическом вакууме.

Порог существенно зависит от объявленной reference normalization. В общем
виде \(\rho_X^{crit}(\rho_{ref})\) определяется уравнением
\begin{equation}
 I_{\rho_X^{crit}}(1/2)
 =3\left[
 8I_{\rho_{ref}}(1/2)-6I_{\rho_{ref}}(1/3)
 \right].
\end{equation}
Для двух используемых benchmark получаем
\begin{align}
 \rho_{ref}=1:
 &\quad \rho_X^{crit}=0.6259781084\ldots,
 &&m_{phys}R_1=0.8011545371\ldots,\\
 \rho_{ref}=\frac12:
 &\quad \rho_X^{crit}=0.1399290937\ldots,
 &&m_{phys}R_1=0.5192110855\ldots.
\end{align}
При уменьшении \(\rho_{ref}\) положительный minimal-content determinant
растёт, поэтому допустимый \(\rho_X\) должен уменьшаться. Значение около
\(0.74\) не удовлетворяет threshold equation для \(\rho_{ref}=1/2\).

\section{Stiffness embedding и точное \(\delta_X=0\)}

Расширенная Hilbert space имеет block form
\begin{equation}
 \mathcal H_{ext}=\mathcal H_{24}\oplus\mathcal H_X.
\end{equation}
Защитная \(\mathbb Z_4^X\) запрещает mixing между blocks. Поэтому neutrino
collective tangent продолжается нулём:
\begin{equation}
 \Xi_{ext}=\Xi\oplus0_X.
\end{equation}
Для canonical product metric
\begin{equation}
 \|\Xi_{ext}\|^2
 =\|\Xi\|^2+\|0_X\|^2
 =23+\pi^{-1},
\end{equation}
то есть
\begin{equation}
 \boxed{\delta_X=0.}
\end{equation}
Тот же вывод следует для determinant Hessian по neutrino collective
coordinate \(t\): при фиксированном выбранном \(\alpha=\pi\)
\begin{equation}
 \frac{\partial^2\Gamma_X}{\partial t^2}=0,
\end{equation}
поскольку \(D_X\) не зависит от \(t\). Если же отождествить \(t\) с Wilson
angle \(\alpha\), появляется \(\rho_X\)-dependent Hessian и denominator
теряет каноническую форму. Поэтому X-sector допустим только как orthogonal
vacuum selector, а \(\pi^{-1}\) по-прежнему приходит из исходной cycle line,
не из X determinant.

\section{Kinetic-mixing и EW/QCD second-sector gate}

Пара \(X_\pm\) является Standard-Model singlet:
\begin{equation}
 (SU(3),SU(2))_Y=(\mathbf1,\mathbf1)_0.
\end{equation}
Поэтому
\begin{equation}
 \Delta\sum Y(B-L)=0,
 \qquad
 \Delta b_1=\Delta b_2=\Delta b_3=0.
\end{equation}
При отсутствии portal couplings её one-loop logarithmic и finite
nonlogarithmic matching vector для Standard-Model gauge couplings равен
\begin{equation}
 \Delta_X^{EW/QCD}=(0,0,0),
\end{equation}
что не может воспроизвести corrected target direction
\begin{equation}
 (1,-1,-6.68221\ldots).
\end{equation}
Двухпетлевая передача через \(B-L\) kinetic mixing возможна только после
задания \(g_{B-L}\), mixing boundary value и portal structure; она не
является parameter-free вторым сектором.

\begin{nogocriterion}[SM-singlet second-sector no-go]
Минимальная vectorlike \(X_\pm\)-пара проходит vacuum-selection и сохраняет
neutrino stiffness, но не даёт независимого EW/QCD normalization-sensitive
следа. Поэтому она остаётся самосогласованным vacuum-selector extension и не
проходит two-sector gate версии III.
\end{nogocriterion}

Здесь необходимы две статусные оговорки. Во-первых, determinant пары выбирает
Wilson branch \(\alpha=\pi\), но не family axis \(n\). Ось по-прежнему должна
быть выведена из residual-orbit dynamics или boundary equations. Во-вторых,
условие на \(\rho_X=M_XR_1\) использует непрерывную массу. Пока \(M_XR_1\) не
получена из parent action, неверно говорить, что селектор не вводит
непрерывных данных: он не требует нового fitting coefficient в stiffness, но
его determinant strength остаётся mass-dependent input.

Следовательно, консолидированный статус имеет вид
\begin{equation}
 \boxed{
 \begin{gathered}
 X_\pm:\ \text{условный динамический selector ветви }\alpha=\pi,\\
 \delta_X=0,\qquad
 \text{family-axis gate открыт},\\
 \text{EW/QCD two-sector gate закрыт отрицательно},\\
 N_{\rm closed\ physical}=0.
 \end{gathered}}
\end{equation}
Положительный результат состоит в существовании anomaly-safe знакового
механизма. Он не является физическим предсказанием и не повышает научный
ledger до вывода \(M_X\), происхождения \(\mathbb Z_4^X\)-line и общего
parent action.

\begin{protocol}[Periodic selector compatibility gate]
Кандидат проходит continuous/discrete anomaly, determinant-sign и
projected-kernel tests. Для физического статуса ещё необходимо:
\begin{enumerate}
 \item вывести \(\mathbb Z_4^X\)-line и её holonomy \(-1\) из finite algebra
 или обязательного topological sector;
 \item вычислить \(M_XR_1\), а не выбирать его ниже threshold;
 \item расширить механизм так, чтобы второй normalization-sensitive sector
 возникал без назначения Standard-Model charges или portals после просмотра
 target.
\end{enumerate}
До этого это минимальная согласованная новая модель, а не восстановление
версии II.A.
\end{protocol}

\section{Collective stiffness на маршруте A}

Пусть \(U_A\) отображает каждый внутренний eigenchannel в соответствующую
нормированную plane wave на внешнем \(S^1\). Для всех spin-, gauge- и
family-shifts
\begin{equation}
 U_A^\dagger U_A=I_{24},
 \qquad
 \|\varphi_{k,\sigma}\|_{L^2(S^1)}=1.
\end{equation}
Поэтому Hilbert--Schmidt traces проекторов не меняются:
\begin{equation}
 \operatorname{Tr}\!\left[(U_AP_{ker}U_A^\dagger)^2\right]=1,
 \qquad
 \operatorname{Tr}\!\left[(U_AP_HU_A^\dagger)^2\right]=23.
\end{equation}
Дополнительный круг даёт единичный множитель нормы и не меняет intrinsic
cycle Gram metric на \(\gamma\):
\begin{equation}
 \|\widehat e_0\|^2=1,
 \qquad
 \|e_1\|^2=L_\gamma^{-1}=\pi^{-1}.
\end{equation}
Следовательно, в canonical product trace--Hodge metric
\begin{equation}
 \boxed{\|\Xi_A\|^2=23+\pi^{-1}.}
\end{equation}
Результат не зависит от \(\theta_*\), spin shift или компенсирующей
\(\mathbb Z_2\)-holonomy: они меняют phases нормированных eigenfunctions, но
не их нормы и не ranks проекторов.

Тем самым collective-stiffness gate маршрута A закрыт как точное
кинематическое тождество canonical metric. Он ещё не возвращает
\(23+\pi^{-1}\) статус физического neutrino denominator: для этого та же
parent action должна выбрать compensated branch, вывести относительную
trace--Hodge metric и пройти второй normalization-sensitive sector.

Второй маршрут сохраняет holonomy на \(\gamma\), но ambient connection
должна иметь кривизну или singular defect source. Если \(2\gamma=\partial
\Sigma_2\) и curvature редуцирована к abelian subgroup с осью \(n\), то
\begin{equation}
 \exp(2\theta_*[n]_\times)
 =\exp\left([n]_\times\int_{\Sigma_2}F_n\right),
\end{equation}
то есть требуется
\begin{equation}
 \int_{\Sigma_2}F_n=2\theta_*\pmod{2\pi}.
\end{equation}
В principal branch это flux
\begin{equation}
 2\theta_*-2\pi=-1.2696746116\ldots.
\end{equation}
Его величина и совпадение оси с factor-selected \(n\) должны следовать из
уравнений движения; иначе это два новых boundary inputs.

Итак, плоский систолический bridge закрыт отрицательно. Положительный
family-selector route теперь имеет две точные версии: плоская связь на
неторсионном \(S^1\), где compensated branch существует и сохраняет
\(23+\pi^{-1}\), а новый \(X_\pm\)-sector может условно выбрать нужную ветвь,
но не выводит axis, mass или второй сектор; либо curved/singular
connection на \(\gamma\) с flux-and-axis dynamical gate. Маршрут A является
алгебраически минимальным, но его минимальный SM-singlet selector закрыт как
two-sector extension. Дальнейшее замыкание требует единого parent action, а
не ещё одного независимого selector-а. Маршрут B остаётся резервным.

Машинный аудит сохранён в
\texttt{s2t\_family\_connection\_defect\_gap\_results.json}.